%=============================================================================
% Chapter 3 Exercises - With hints from original book (30 exercises)
%=============================================================================
\section*{Exercises}
\addcontentsline{toc}{section}{Exercises}

\begin{enumerate}
    \item \label{exer1-chapter3}Let $S$ be a multiplicatively closed subset of a ring $A$, and let $M$ be a finitely generated $A$-module. Prove that $S^{-1}M = 0$ if and only if there exists $s \in S$ such that $sM = 0$.

    \item \label{exer2-chapter3}Let $\mathfrak{a}$ be an ideal of a ring $A$, and let $S = 1 + \mathfrak{a}$. Show that $S^{-1}\mathfrak{a}$ is contained in the Jacobson radical of $S^{-1}A$.\\
    \textit{Hint:} Use this result and Nakayama's lemma to give a proof of \eqref{cor:cayley-hamilton} which does not depend on determinants. [If $M = \mathfrak{a}M$, then $S^{-1}M = (S^{-1}\mathfrak{a})(S^{-1}M)$, hence by Nakayama we have $S^{-1}M = 0$. Now use Exercise \ref{exer1-chapter3}.]

    \item \label{exer3-chapter3}Let $A$ be a ring, let $S$ and $T$ be two multiplicatively closed subsets of $A$, and let $U$ be the image of $T$ in $S^{-1}A$. Show that the rings $(ST)^{-1}A$ and $U^{-1}(S^{-1}A)$ are isomorphic.

    \item \label{exer4-chapter3}Let $f \colon A \to B$ be a homomorphism of rings and let $S$ be a multiplicatively closed subset of $A$. Let $T = f(S)$. Show that $S^{-1}B$ and $T^{-1}B$ are isomorphic as $S^{-1}A$-modules.

    \item \label{exer5-chapter3}Let $A$ be a ring. Suppose that, for each prime ideal $\mathfrak{p}$, the local ring $A_{\mathfrak{p}}$ has no nilpotent element $\neq 0$. Show that $A$ has no nilpotent element $\neq 0$. If each $A_{\mathfrak{p}}$ is an integral domain, is $A$ necessarily an integral domain?

    \item \label{exer6-chapter3}Let $A$ be a ring $\neq 0$ and let $\Sigma$ be the set of all multiplicatively closed subsets $S$ of $A$ such that $0 \notin S$. Show that $\Sigma$ has maximal elements, and that $S \in \Sigma$ is maximal if and only if $A \setminus S$ is a minimal prime ideal of $A$".

    \item \label{exer7-chapter3}A multiplicatively closed subset $S$ of a ring $A$ is said to be \emph{saturated}\index{Saturated} if $xy \in S \iff x \in S \text{ and } y \in S$. Prove that:
    \begin{enumerate}[i)]
        \item $S$ is saturated $\iff$ $A \setminus S$ is a union of prime ideals.
        \item If $S$ is any multiplicatively closed subset of $A$, there is a unique smallest saturated multiplicatively closed subset $\bar{S}$ containing $S$, and that $\bar{S}$ is the complement in $A$ of the union of the prime ideals which do not meet $S$. ($\bar{S}$ is called the \emph{saturation} of $S$.)
    \end{enumerate}
    If $S = 1 + \mathfrak{a}$, where $\mathfrak{a}$ is an ideal of $A$, find $\bar{S}$.

    \item \label{exer8-chapter3} Let $S, T$ be multiplicatively closed subsets of $A$, such that $S \subseteq T$. Let $\phi \colon S^{-1}A \to T^{-1}A$ be the homomorphism which maps each $a/s \in S^{-1}A$ to $a/s$ considered as an element of $T^{-1}A$. Show that the following statements are equivalent:
    \begin{enumerate}[i)]
        \item $\phi$ is bijective.
        \item For each $t \in T$, $t/1$ is a unit in $S^{-1}A$.
        \item For each $t \in T$ there exists $x \in A$ such that $xt \in S$.
        \item $T$ is contained in the saturation of $S$ (Exercise \ref{exer7-chapter3}).
        \item Every prime ideal which meets $T$ also meets $S$.
    \end{enumerate}

    \item \label{exer9-chapter3} The set $S_0$ of all non-zero-divisors in $A$ is a saturated multiplicatively closed subset of $A$. Hence the set $D$ of zero-divisors in $A$ is a union of prime ideals (see Chapter \ref{chapter1}, Exercise \ref{exer14-chapter1}). Show that every minimal prime ideal of $A$ is contained in $D$. \textit{Hint}: Use Exercise \ref{exer6-chapter3}.
    
    \quad\, The ring $S_0^{-1}A$ is called the \emph{total ring of fractions} of $A$. Prove that:
    \begin{enumerate}[i)]
        \item $S_0$ is the largest multiplicatively closed subset of $A$ for which the homomorphism $A \to S_0^{-1}A$ is injective.
        \item Every element in $S_0^{-1}A$ is either a zero-divisor or a unit.
        \item Every ring in which every non-unit is a zero-divisor is equal to its total ring of fractions (i.e., $A \to S_0^{-1}A$ is bijective).
    \end{enumerate}

    \item \label{exer10-chapter3} Let $A$ be a ring.
    \begin{enumerate}[i)]
        \item If $A$ is absolutely flat (Chapter \ref{chapter2}, Exercise \ref{exer27-chapter2}) and $S$ is any multiplicatively closed subset of $A$, then $S^{-1}A$ is absolutely flat.
        \item $A$ is absolutely flat $\iff$ $A_{\mathfrak{m}}$ is a field for each maximal ideal $\mathfrak{m}$.
    \end{enumerate}

    \item \label{exer11-chapter3} Let $A$ be a ring. Prove that the following are equivalent:
    \begin{enumerate}[i)]
        \item $A/\mathfrak{N}$ is absolutely flat ($\mathfrak{N}$ being the nilradical of $A$).
        \item Every prime ideal of $A$ is maximal.
        \item $\Spec(A)$ is a $T_1$-space (i.e., every subset consisting of a single point is closed).
        \item $\Spec(A)$ is Hausdorff.
    \end{enumerate}
    If these conditions are satisfied, show that $\Spec(A)$ is compact and totally disconnected (i.e., the only connected subsets of $\Spec(A)$ are those consisting of a single point).

    \item \label{exer12-chapter3} Let $A$ be an integral domain and $M$ an $A$-module. An element $x \in M$ is a \emph{torsion element}\index{Torsion, element} of $M$ if $\Ann(x) \neq 0$, that is if $x$ is killed by some non-zero element of $A$. Show that the torsion elements of $M$ form a submodule of $M$. This submodule is called the \emph{torsion submodule}\index{Torsion, element!submodule} of $M$ and is denoted by $T(M)$. If $T(M) = 0$, the module $M$ is said to be \emph{torsion-free}. Show that:
    \begin{enumerate}[i)]
        \item If $M$ is any $A$-module, then $M/T(M)$ is torsion-free.
        \item If $f \colon M \to N$ is a module homomorphism, then $f(T(M)) \subseteq T(N)$.
        \item If $0 \to M' \to M \to M''$ is an exact sequence, then the sequence $0 \to T(M') \to T(M) \to T(M'')$ is exact.
        \item If $M$ is any $A$-module, then $T(M)$ is the kernel of the mapping $x \mapsto 1 \otimes x$ of $M$ into $K \otimes_A M$, where $K$ is the field of fractions of $A$.
        
        \textit{Hint:} For iv), show that $K$ may be regarded as the direct limit of its submodules $Ag$ ($g \in K$); using Chapter \ref{chapter1}, Exercise \ref{exer15-chapter1} and Exercise \ref{exer20-chapter1}, show that if $1 \otimes x = 0$ in $K \otimes M$ then $1 \otimes x = 0$ in $Ag \otimes M$ for some $g \neq 0$. Deduce that $g^{-1}x = 0$.
    \end{enumerate}

    \item \label{exer13-chapter3} Let $S$ be a multiplicatively closed subset of an integral domain $A$. In the notation of Exercise \ref{exer12-chapter3}, show that $T(S^{-1}M) = S^{-1}(T(M))$. Deduce that the following are equivalent:
    \begin{enumerate}[i)]
        \item $M$ is torsion-free.
        \item $M_{\mathfrak{p}}$ is torsion-free for all prime ideals $\mathfrak{p}$.
        \item $M_{\mathfrak{m}}$ is torsion-free for all maximal ideals $\mathfrak{m}$.
    \end{enumerate}

    \item \label{exer14-chapter3} Let $M$ be an $A$-module and $\mathfrak{a}$ an ideal of $A$. Suppose that $M_{\mathfrak{m}} = 0$ for all maximal ideals $\mathfrak{m} \not\supseteq \mathfrak{a}$. Prove that $M = \mathfrak{a}M$.
    
    \textit{Hint:} Pass to the $A/\mathfrak{a}$-module $M/\mathfrak{a}M$ and use \eqref{prop:zero-local}.

    \item \label{exer15-chapter3}Let $A$ be a ring, and let $F$ be the $A$-module $A^n$. Show that every set of $n$ generators of $F$ is a basis of $F$.
    
    \textit{Hint:} Let $x_1, \ldots, x_n$ be a set of generators and $e_1, \ldots, e_n$ the canonical basis of $F$. Define $\phi \colon F \to F$ by $\phi(e_i) = x_i$. Then $\phi$ is surjective and we have to prove that it is an isomorphism. By \eqref{prop:injective-surjective-local} we may assume that $A$ is a local ring. Let $N$ be the kernel of $\phi$ and let $k = A/\mathfrak{m}$ be the residue field of $A$. Since $F$ is a flat $A$-module, the exact sequence $0 \to N \to F \xrightarrow{\phi} F \to 0$ gives an exact sequence $0 \longrightarrow k \otimes N \longrightarrow k \otimes F \xrightarrow{1 \otimes \phi} k \otimes F \longrightarrow 0$. Now $k \otimes F = k^n$ is an $n$-dimensional vector space over $k$; $1 \otimes \phi$ is surjective, hence bijective, hence $k \otimes N = 0$. Also $N$ is finitely generated, by Chapter \ref{chapter2}, Exercise \ref{exer12-chapter2}, hence $N = 0$ by Nakayama's lemma. Hence $\phi$ is an isomorphism.

    Deduce that every set of generators of $F$ has at least $n$ elements.

    \item \label{exer16-chapter3}Let $B$ be a flat $A$-algebra. Then the following conditions are equivalent:
    \begin{enumerate}[i)]
        \item $\mathfrak{a}^{ec} = \mathfrak{a}$ for all ideals $\mathfrak{a}$ of $A$.
        \item $\Spec(B) \to \Spec(A)$ is surjective.
        \item For every maximal ideal $\mathfrak{m}$ of $A$ we have $\mathfrak{m}^e \neq (1)$.
        \item If $M$ is any non-zero $A$-module, then $M_B \neq 0$.
        \item For every $A$-module $M$, the mapping $x \mapsto 1 \otimes x$ of $M$ into $M_B$ is injective.
    \end{enumerate}
    \textit{Hint:} For i) $\implies$ ii), use \eqref{prop:extension-ideals}. 
    
    ii) $\implies$ iii) is clear. 
    
    iii) $\implies$ iv): Let $x$ be a non-zero element of $M$ and let $M' = Ax$. Since $B$ is flat over $A$ it is enough to show that $M'_B \neq 0$. We have $M' \cong A/\mathfrak{a}$ for some ideal $\mathfrak{a} \neq (1)$, hence $M'_B \cong B/\mathfrak{a}^e$. Now $\mathfrak{a} \subseteq \mathfrak{m}$ for some maximal ideal $\mathfrak{m}$, hence $\mathfrak{a}^e \subseteq \mathfrak{m}^e \neq (1)$. Hence $M'_B \neq 0$. 
    
    iv) $\implies$ v): Let $M'$ be the kernel of $M \to M_B$. Since $B$ is flat over $A$, the sequence $0 \to M'_B \to M_B \to (M_B)_B$ is exact. But (Chapter \ref{chapter2}, Exercise \ref{exer13-chapter2}, with $N = M_B$) the mapping $M_B \to (M_B)_B$ is injective, hence $M'_B = 0$ and therefore $M' = 0$. 
    
    v) $\implies$ i): Take $M = A/\mathfrak{a}$.

    $B$ is said to be \emph{faithfully flat}\index{Flat!faithfully} over $A$.

    \item \label{exer17-chapter3}Let $A \xrightarrow{f} B \xrightarrow{g} C$ be ring homomorphisms. If $g \circ f$ is flat and $g$ is faithfully flat, then $f$ is flat.

    \item \label{exer18-chapter3}Let $f \colon A \to B$ be a flat homomorphism of rings, let $\mathfrak{q}$ be a prime ideal of $B$ and let $\mathfrak{p} = \mathfrak{q}^c$. Then $f^* \colon \Spec(B_{\mathfrak{q}}) \to \Spec(A_{\mathfrak{p}})$ is surjective.\\
    \textit{Hint:} For $B_{\mathfrak{p}}$ is flat over $A_{\mathfrak{p}}$ by \eqref{prop:flat-local}, and $B_{\mathfrak{q}}$ is a local ring of $B_{\mathfrak{p}}$, hence is flat over $B_{\mathfrak{p}}$. Hence $B_{\mathfrak{q}}$ is flat over $A_{\mathfrak{p}}$ and satisfies condition (3) of Exercise \ref{exer16-chapter3}.

    \item \label{exer19-chapter3}Let $A$ be a ring, $M$ an $A$-module. The \emph{support}\index{Support} of $M$ is defined to be the set $\Supp(M)$ of prime ideals $\mathfrak{p}$ of $A$ such that $M_{\mathfrak{p}} \neq 0$. Prove the following results:
    \begin{enumerate}[i)]
        \item $M \neq 0 \iff \Supp(M) \neq \varnothing$.
        \item $V(\mathfrak{a}) = \Supp(A/\mathfrak{a})$.
        \item If $0 \to M' \to M \to M'' \to 0$ is an exact sequence, then $\Supp(M) = \Supp(M') \cup \Supp(M'')$.
        \item If $M = \bigoplus M_i$, then $\Supp(M) = \bigcup \Supp(M_i)$.
        \item If $M$ is finitely generated, then $\Supp(M) = V(\Ann(M))$ (and is therefore a closed subset of $\Spec(A)$).
        \item If $M, N$ are finitely generated, then $\Supp(M \otimes_A N) = \Supp(M) \cap \Supp(N)$. [Use Chapter \ref{chapter2}, Exercise \ref{exer3-chapter2}.]
        \item If $M$ is finitely generated and $\mathfrak{a}$ is an ideal of $A$, then $\Supp(M/\mathfrak{a}M) = V(\mathfrak{a} + \Ann(M))$.
        \item If $f \colon A \to B$ is a ring homomorphism and $M$ is a finitely generated $A$-module, then $\Supp(B \otimes_A M) = f^{*-1}(\Supp(M))$.
    \end{enumerate}

    \item \label{exer20-chapter3}Let $f \colon A \to B$ be a ring homomorphism, $f^* \colon \Spec(B) \to \Spec(A)$ the associated mapping. Show that:
    \begin{enumerate}[i)]
        \item Every prime ideal of $A$ is a contracted ideal $\iff$ $f^*$ is surjective.
        \item Every prime ideal of $B$ is an extended ideal $\implies$ $f^*$ is injective.
    \end{enumerate}
    Is the converse of ii) true?

    \item \label{exer21-chapter3} \begin{enumerate}[i)]
        \item Let $A$ be a ring, $S$ a multiplicatively closed subset of $A$, and $\phi \colon A \to S^{-1}A$ the canonical homomorphism. Show that $\phi^* \colon \Spec(S^{-1}A) \to \Spec(A)$ is a homeomorphism of $\Spec(S^{-1}A)$ onto its image in $X = \Spec(A)$. Let this image be denoted by $S^{-1}X$. In particular, if $f \in A$, the image of $\Spec(A_f)$ in $X$ is the basic open set $X_f$ (Chapter \ref{chapter1}, Exercise \ref{exer17-chapter1}).
        \item Let $f \colon A \to B$ be a ring homomorphism. Let $X = \Spec(A)$ and $Y = \Spec(B)$, and let $f^* \colon Y \to X$ be the mapping associated with $f$. Identifying $\Spec(S^{-1}A)$ with its canonical image $S^{-1}X$ in $X$, and $\Spec(S^{-1}B)$ ($= \Spec(f(S)^{-1}B)$) with its canonical image $S^{-1}Y$ in $Y$, show that $S^{-1}f^* \colon \Spec(S^{-1}B) \to \Spec(S^{-1}A)$ is the restriction of $f^*$ to $S^{-1}Y$, and that $S^{-1}Y = f^{*-1}(S^{-1}X)$.
        \item Let $\mathfrak{a}$ be an ideal of $A$ and let $\mathfrak{b} = \mathfrak{a}^e$ be its extension in $B$. Let $\bar{f} \colon A/\mathfrak{a} \to B/\mathfrak{b}$ be the homomorphism induced by $f$. If $\Spec(A/\mathfrak{a})$ is identified with its canonical image $V(\mathfrak{a})$ in $X$, and $\Spec(B/\mathfrak{b})$ with its image $V(\mathfrak{b})$ in $Y$, show that $\bar{f}^*$ is the restriction of $f^*$ to $V(\mathfrak{b})$.
        \item Let $\mathfrak{p}$ be a prime ideal of $A$. Take $S = A \setminus \mathfrak{p}$ in ii) and then reduce mod $S^{-1}\mathfrak{p}$ as in iii). Deduce that the subspace $f^{*-1}(\mathfrak{p})$ of $Y$ is naturally homeomorphic to $\Spec(B_{\mathfrak{p}}/\mathfrak{p}B_{\mathfrak{p}}) = \Spec(k(\mathfrak{p}) \otimes_A B)$, where $k(\mathfrak{p})$ is the residue field of the local ring $A_{\mathfrak{p}}$. $\Spec(k(\mathfrak{p}) \otimes_A B)$ is called the \emph{fiber} of $f^*$ over $\mathfrak{p}$.
    \end{enumerate}

    \item \label{exer22-chapter3}Let $A$ be a ring and $\mathfrak{p}$ a prime ideal of $A$. Then the canonical image of $\Spec(A_{\mathfrak{p}})$ in $\Spec(A)$ is equal to the intersection of all the open neighborhoods of $\mathfrak{p}$ in $\Spec(A)$.

    \item \label{exer23-chapter3}Let $A$ be a ring, let $X = \Spec(A)$ and let $U$ be a basic open set in $X$ (i.e., $U = X_f$ for some $f \in A$: Chapter \ref{chapter1}, Exercise \ref{exer17-chapter1}).
    \begin{enumerate}[i)]
        \item If $U = X_f$, show that the ring $A(U) = A_f$ depends only on $U$ and not on $f$.
        \item Let $U' = X_g$ be another basic open set such that $U' \subseteq U$. Show that there is an equation of the form $g^n = uf$ for some integer $n > 0$ and some $u \in A$, and use this to define a homomorphism $\rho \colon A(U) \to A(U')$ (i.e., $A_f \to A_g$) by mapping $a/f^m$ to $au^m/g^{mn}$. Show that $\rho$ depends only on $U$ and $U'$. This homomorphism is called the \emph{restriction homomorphism}.
        \item If $U = U'$, then $\rho$ is the identity map.
        \item If $U \supseteq U' \supseteq U''$ are basic open sets in $X$, show that the diagram (with arrows being restriction homomorphisms) is commutative.
        \[\begin{tikzcd}
            {A(U)} && {A(U'')} \\
            & {A(U')}
            \arrow[from=1-1, to=1-3]
            \arrow[from=1-1, to=2-2]
            \arrow[from=2-2, to=1-3]
        \end{tikzcd}\]
        \item Let $x$ ($= \mathfrak{p}$) be a point of $X$. Show that \[\varinjlim_{U\ni x} A(U) \cong A_{\mathfrak{p}}.\]
    \end{enumerate}
    The assignment of the ring $A(U)$ to each basic open set $U$ of $X$, and the restriction homomorphisms $\rho$, satisfying the conditions iii) and iv) above, constitutes a \emph{presheaf of rings} on the basis of open sets $(X_f)_{f \in A}$. v) says that the \emph{stalk} of this presheaf at $x \in X$ is the corresponding local ring $A_{\mathfrak{p}}$.

    \item \label{exer24-chapter3}Show that the presheaf of Exercise \ref{exer23-chapter3} has the following property. Let $(U_i)_{i \in I}$ be a covering of $X$ by basic open sets. For each $i \in I$ let $s_i \in A(U_i)$ be such that, for each pair of indices $i,j$, the images of $s_i$ and $s_j$ in $A(U_i \cup U_j)$ are equal. Then there exists a unique $s \in A$ ($= A(X)$) whose image in $A(U_i)$ is $s_i$ for all $i \in I$. (This essentially implies that the presheaf is a \emph{sheaf}.)

    \item \label{exer25-chapter3}Let $f \colon A \to B$, $g \colon A \to C$ be ring homomorphisms and let $h \colon A \to B \otimes_A C$ be defined by $h(x) = f(x) \otimes g(x)$. Let $X, Y, Z, T$ be the prime spectra of $A, B, C, B \otimes_A C$ respectively. Then $h^*(T) = f^*(Y) \cap g^*(Z)$.
    
    \textit{Hint:} Let $\mathfrak{p} \in X$, and let $k = k(\mathfrak{p})$ be the residue field at $\mathfrak{p}$. By Exercise 21, the fiber $h^{*-1}(\mathfrak{p})$ is the spectrum of $(B \otimes_A C) \otimes_A k \cong (B \otimes_A k) \otimes_k (C \otimes_A k)$. Hence $\mathfrak{p} \in h^*(T) \iff (B \otimes_A k) \otimes_k (C \otimes_A k) \neq 0 \iff B \otimes_A k \neq 0 \text{ and } C \otimes_A k \neq 0 \iff \mathfrak{p} \in f^*(Y) \cap g^*(Z)$.

    \item \label{exer26-chapter3}Let $(B_\alpha, g_{\alpha\beta})$ be a direct system of rings and $B$ the direct limit. For each $\alpha$, let $f_\alpha \colon A \to B_\alpha$ be a ring homomorphism such that $g_{\alpha\beta} \circ f_\alpha = f_\beta$ whenever $\alpha \leq \beta$ (i.e., the $B_\alpha$ form a direct system of $A$-algebras). The $f_\alpha$ induce $f \colon A \to B$. Show that
    \[
    f^*(\Spec(B)) = \bigcap_\alpha f_\alpha^*(\Spec(B_\alpha)).
    \]
    \textit{Hint:} Let $\mathfrak{p} \in \Spec(A)$. Then $f^{*-1}(\mathfrak{p})$ is the spectrum of \[B \otimes_A k(\mathfrak{p}) \cong \varinjlim (B_\alpha \otimes_A k(\mathfrak{p}))\] (since tensor products commute with direct limits: Chapter \ref{chapter2}, Exercise \ref{exer20-chapter2}). By Exercise \ref{exer21-chapter2} of Chapter \ref{chapter2} it follows that $f^{*-1}(\mathfrak{p}) = \varnothing$ if and only if $B_\alpha \otimes_A k(\mathfrak{p}) = 0$ for some $\alpha$, i.e., if and only if $f_\alpha^{*-1}(\mathfrak{p}) = \varnothing$.

    \item \label{exer27-chapter3} \begin{enumerate}[i)]
        \item Let $f_\alpha \colon A \to B_\alpha$ be any family of $A$-algebras and let $f \colon A \to B$ be their tensor product over $A$ (Chapter \ref{chapter2}, Exercise \ref{exer23-chapter2}). Then \[f^*(\Spec(B)) = \bigcap_\alpha f_\alpha^*(\Spec(B_\alpha)).\] [Use Exercises \ref{exer25-chapter3} and \ref{exer26-chapter3}.]
        \item Let $f_\alpha \colon A \to B_\alpha$ be any finite family of $A$-algebras and let $B = \prod_\alpha B_\alpha$. Define $f \colon A \to B$ by $f(x) = (f_\alpha(x))$. Then $f^*(\Spec(B)) = \bigcup_\alpha f_\alpha^*(\Spec(B_\alpha))$.
        \item Hence the subsets of $X = \Spec(A)$ of the form $f^*(\Spec(B))$, where $f \colon A \to B$ is a ring homomorphism, satisfy the axioms for closed sets in a topological space. The associated topology is the \emph{constructible topology}\index{Constructible topology} on $X$. It is finer than the Zariski topology (i.e., there are more open sets, or equivalently more closed sets).
        \item Let $X_c$ denote the set $X$ endowed with the constructible topology. Show that $X_c$ is quasi-compact.
    \end{enumerate}

    \item \label{exer28-chapter3} (Continuation of Exercise \ref{exer27-chapter3}.)
    \begin{enumerate}[i)]
        \item For each $g \in A$, the set $X_g$ (Chapter \ref{chapter1}, Exercise \ref{exer17-chapter1}) is both open and closed in the constructible topology.
        \item Let $C'$ denote the smallest topology on $X$ for which the sets $X_g$ are both open and closed, and let $X_{C'}$ denote the set $X$ endowed with this topology. Show that $X_{C'}$ is Hausdorff.
        \item Deduce that the identity mapping $X_c \to X_{C'}$ is a homeomorphism. Hence a subset $E$ of $X$ is of the form $f^*(\Spec(B))$ for some $f \colon A \to B$ if and only if it is closed in the topology $C'$.
        \item The topological space $X_c$ is compact, Hausdorff and totally disconnected.
    \end{enumerate}

    \item Let $f \colon A \to B$ be a ring homomorphism. Show that $f^* \colon \Spec(B) \to \Spec(A)$ is a continuous closed mapping (i.e., maps closed sets to closed sets) for the constructible topology.

    \item Show that the Zariski topology and the constructible topology on $\Spec(A)$ are the same if and only if $A/\mathfrak{N}$ is absolutely flat (where $\mathfrak{N}$ is the nilradical of $A$).\\
    \textit{Hint:} Use Exercise \ref{exer11-chapter3}.
\end{enumerate}
